\section{Introduction}\label{sec:introduction}

Firebreaks interrupt the spread of fire through a landscape, but their
effect depends on where a future fire starts and which propagation
routes it follows. Placement must therefore be decided before the fire
realization is known. Fuel-treatment planning has used spatial
optimization and two-stage stochastic models for uncertain wildfire
behavior \citep{agee2000use,kabli2015stochastic}. A sampled scenario
can preserve the geography of a simulated fire rather than reducing
risk to independent cell scores.

Cell2Fire generates spatially explicit fire realizations that can
support planning decisions \citep{pais2021cell2fire}; downstream network
value also helps identify potentially useful treatment locations
\citep{pais2021downstream}. Multiple routes in a directed propagation
graph matter when assessing a proposed firebreak: blocking one route
need not protect a cell reachable through another. Retaining those
routes across hundreds of scenarios, however, can make optimization
costly. Reduced propagation graphs provide a tractable representation
whose solver behavior and decision quality require separate evaluation.

Here the optimization input labeled \texttt{shortest\_path} is the
\emph{reduced shortest-path graph}; a separate Reburn input denotes the
original propagation representation. The labels do
not make the reduced graph a rooted tree. In the 20-by-20 campaign,
preflight classifies 86.83\% of reduced scenarios as acyclic DAGs that
are not trees, 12.41\% as rooted arborescences and 0.76\% as empty or
invalid. If the graphs in a pair have identical cells, ignition and
weights and the reduced arc set is included in the Reburn arc set, a
fixed placement has no greater loss on the reduced graph. The original
scenario inputs needed to check those premises and evaluate the same
placement on both graphs are absent from the committed results. We
therefore do not infer the quality of the graph reduction from
independently solved reduced and Reburn objectives.

The sampled placement problem remains a mixed-integer program: one
treatment plan must work across many graph realizations and, under
CVaR, limit adverse losses. Branch-and-Benders separates the shared
placement from scenario propagation \citep{rahmaniani2017benders}.
Root relaxations can underestimate losses when recourse cuts are
incomplete. We study structural lower bounds based on downstream
reachability, capped coverage and ignition-to-cell paths, as well as a
combinatorial Benders comparator motivated by network-spread methods
\citep{taninmics2024benders}. The complete path family describes
fractional recourse on a general directed graph. Only on a rooted tree
does complete coverage also coincide with it; practical path
separation retains only a limited set of paths.

This article contributes a stochastic reachability formulation with
relatively complete recourse, proofs of the validity and projection of
the structural bounds, and a numerically audited comparison of exact
methods on matched scenario samples. The complete 810-run reduced
\texttt{new20x20} panel covers three objectives, three training sizes
and six methods per objective. It supplies certification counts,
mean and maximum final gaps, and observed solver-time means and
sample standard deviations, with an explicit budget-level analysis. A
separate 180-run paired sensitivity experiment tests ascending and
descending ordering of current scenario recourse estimates in
Combinatorial B\&B. Its effect is objective-dependent: descending order
improves mean--CVaR certification in this experiment, but not CVaR. A
distinct 1,440-run \texttt{Sub20} panel tests polynomial and separated
projected inequalities. These data support statements about
computational certification on the reported panels; corrupted
placement/evaluation fields, incomplete Reburn runs and missing source
graphs preclude conclusions about out-of-sample placement quality or
the reduced/Reburn loss difference.

Section~\ref{sec:setting} specifies the graphs and conditional reduction
result; Section~\ref{sec:model} develops the stochastic model and
recourse. Section~\ref{sec:methods} proves the bound relationships.
Sections~\ref{sec:computational} and \ref{sec:results} document the
recorded experiments and numerical evidence, followed by conclusions
in Section~\ref{sec:conclusions}.
